\section{绪论}

\subsection{研究背景与问题}

视觉语言模型在近年来取得了前所未有的进展。以 CLIP 为代表的对比式模型通过大规模图文对齐学习，将视觉和语言映射到统一的语义嵌入空间，实现了强大的零样本泛化能力。在此基础上，BLIP-2、LLaVA 和 Qwen-VL 等生成式视觉语言模型进一步实现了从视觉输入到自然语言输出的端到端推理。然而，这些模型的性能高度依赖输入图像的感知质量。在低分辨率、模糊、噪声和遮挡等视觉退化条件下，视觉编码器提取的特征会发生严重畸变，导致下游的语义理解产生系统性偏差。视觉退化条件下的鲁棒多模态理解，是视觉语言模型从实验室走向实际部署需要解决的核心挑战之一。

脑电图以毫秒级时间分辨率记录视觉刺激引发的大脑皮层电活动，已在视觉语义解码领域显示出潜力。已有研究证明，从人类观看自然图像时的 EEG 信号中可以解码物体类别、检索语义相关图像，并与预训练视觉语言模型（如 CLIP）的嵌入空间建立映射。这引出了一个具有理论和应用价值的问题：在视觉输入质量受损时，同步采集的 EEG 信号能否作为辅助信息源，改善视觉语言模型的语义理解性能？

本文围绕上述问题，在小样本 EEG-image 配对数据条件下开展了系统性研究。本文的研究定位是将 EEG 作为一种视觉语义增强信号，通过严格的 shuffled EEG 和 random EEG 对照组设计，为 EEG 辅助 VLM 的有效性提供可验证的经验证据。

\subsection{研究意义}

本文的研究意义体现在两个维度。在学术层面，`EEG 是否携带可被深度学习模型利用的视觉语义信息''是计算神经科学与多模态学习交叉领域的关键问题。现有研究大多在完整视觉条件下进行，未系统考察视觉退化场景。本文通过在六种退化条件下设置四种 EEG 对照模式的严格实验，为该问题提供了结构化的定量证据。在技术层面，本文提出的 A2 temporal-spectral-spatial EEG 编码器提供了 EEG 多维度特征提取的可行方案，VTF token-level 融合方法为将 EEG 信息集成到原生视觉语言模型 pipeline 中提供了技术路径。

\subsection{主要贡献}

本文的主要贡献包括：

\begin{enumerate}
  \item 构建了视觉退化条件下 EEG 辅助语义识别与生成式描述的完整实验框架，包含六种退化条件、四种对照模式和 image-level 数据划分策略；
  \item 设计了 A2 temporal-spectral-spatial EEG 编码器，通过三支路并行编码和后融合 Transformer 实现 EEG 三维特征的端到端提取与交互；
  \item 提出了 VTF token-level EEG 融合方法，使 EEG 信息可在 CLIP ViT visual token 层面与视觉表示进行跨模态交互；
  \item 实现了基于 Qwen2-VL 的生成式视觉语言模型，验证了 EEG-enhanced visual tokens 驱动自由形式 caption 生成的可行性。
\end{enumerate}

\subsection{本文结构}

本文共分十章。第 2 章梳理相关技术基础；第 3 章介绍数据集与任务设定；第 4 章阐述系统总体架构；第 5 章详述 A2 EEG 语义融合模型；第 6 章介绍 token-level VTF 融合方法；第 7 章描述生成式 VLM 的设计与多种技术方案对比；第 8 章报告完整实验结果与分析；第 9 章概述工程实现；第 10 章总结全文并提出未来方向。
